\section{Введение}
\label{sec:introduction}

Устойчивый интерес к полупроводниковым квантовым точкам (КТ) связан
с дискретностью их электронных состояний и возможностью изменять
оптические свойства выбором состава и размера нанокристалла. Подобно
искусственному атому, КТ обладает резонансными переходами, однако допускает
химическую функционализацию и включение в твёрдую матрицу или составную
наноструктуру. Для управления её оптическим откликом существенно не только
положение экситонного перехода, но и поле, непосредственно действующее
на этот переход. Плазмонные металлические наночастицы (МНЧ) позволяют
концентрировать свет в ближней зоне с размерами существенно меньше длины
волны и тем самым изменять эффективность возбуждения расположенной рядом КТ.

Объединение КТ и МНЧ создаёт связанную экситон-плазмонную систему,
оптический отклик которой не сводится к сумме откликов компонентов.
Взаимодействие узкого экситонного перехода с более широким плазмонным
резонансом может приводить к сдвигу и уширению линии, интерференционным
особенностям типа Фано и нелинейному изменению спектра при насыщении КТ.
Такие эффекты были рассмотрены, в частности, в работе Zhang, Govorov
и Bryant~\cite{zhang2006}. Экспериментальные исследования химически
собранных комплексов КТ--золото показали возможность раздельно выявлять
поляризационно-зависимое усиление возбуждающего поля и изменение времени
распада экситона~\cite{cohenhoshen2012}. Эти два проявления связи нельзя
отождествлять: усиление возбуждения может сопровождаться дополнительными
потерями энергии и уменьшением квантового выхода. Поэтому увеличение
яркости люминесценции, снижение порога создания экситона и достижение режима
сильной связи представляют собой разные физические задачи.

Особенно важен переход от стационарного возбуждения к воздействию
ультракороткого импульса. В работе Shah и соавторов~\cite{shah2013}
сопоставлены квантовая и полуклассическая модели КТ, связанной с плазмонной
наноантенной, и предсказано обращение резонанса Фано при фемтосекундном
возбуждении: провал в оптическом спектре сменяется пиком вследствие изменения
фазового соотношения дипольных откликов. Рассмотренная авторами структура
содержала КТ между двумя золотыми вытянутыми сфероидами. В полуклассическом
описании отклик всей металлической антенны представлялся одним эффективным
затухающим осциллятором, параметры которого определялись по электродинамическому
расчёту. Авторы также показали, что изменение скорости спонтанного распада
вблизи металла существенно для сопоставления возбуждения и релаксации.
Таким образом, эта работа задаёт как основу для описания импульсной
динамики, так и условия осторожного переноса её выводов на другие геометрии.

Для композитов из коллоидных КТ и золотых наноцилиндров практически важен
более простой элемент: одна КТ вблизи одной МНЧ. Возникает вопрос, где следует
разместить КТ и как выбрать поляризацию лазера, чтобы получить заданное
возбуждение при меньшем падающем флюенсе. Здесь цилиндр заменяется
осесимметричным вытянутым эллипсоидом -- сфероидом. Такая идеализация
сохраняет различие продольного и поперечного плазмонных откликов и позволяет
последовательно уточнять пространственное описание поля. Она не устраняет
различие между реальной формой цилиндра и сфероидом, но даёт контролируемое
сравнение приближений при одной и той же заданной поверхности.

В диполь-дипольной модели МНЧ заменяется диполем в её центре, а связь
определяется расстоянием между центрами и ориентацией дипольных моментов.
Для КТ в непосредственной близости от поверхности этой замены может быть
недостаточно. В расширенной модели используется полное квазистатическое
решение для сфероида: учитываются пространственная неоднородность его поля
и высшие мультипольные составляющие реакции на диполь КТ. Сама КТ в обоих
случаях остаётся точечным диполем, а поле вычисляется в её центре.
Следовательно, уточняется именно пространственная модель металлической
частицы; конечный размер КТ и запаздывание требуют дополнительных
приближений и здесь не включаются.

Независимое уточнение относится к дисперсии металла. Короткий импульс
возбуждает конечную полосу частот, поэтому важны не только значение
поляризуемости на несущей, но и её комплексная частотная зависимость.
Мы используем табличные оптические константы золота
Джонсона--Кристи~\cite{johnson1972} и представляем поляризуемость причинной
пассивной суммой осцилляторов. Это позволяет сравнить одноосцилляторное
и многоосцилляторное описания при неизменной геометрии. Число осцилляторов
характеризует точность представления дисперсии и не совпадает с числом
пространственных мультиполей. Сопоставление относится к заданной полосе
подгонки и не означает, что один эффективно подобранный осциллятор
непригоден для любого узкополосного эксперимента.

Цель настоящей работы -- определить, как взаимное расположение КТ и золотой
МНЧ, поляризация поля и точность описания дисперсии влияют на резонансное
нелинейное возбуждение КТ коротким импульсом. Рассчитываются динамика
населённости, порог по флюенсу, слабополевые спектральные признаки
и оптическая работа падающего поля. Для выбранного набора параметров
устанавливается, когда дипольное приближение воспроизводит полное
квазистатическое решение, и формулируются рекомендации по размещению КТ
и ориентации МНЧ. Практическая часть посвящена спектроскопическому
различению преимущественных конфигураций в ансамбле и способам их получения
химической сборкой. В расчёте задаётся направление электрического поля,
а не волнового вектора; выбор направления освещения экспериментального
слоя обсуждается отдельно как способ реализовать требуемую поляризацию.
